%% ma: L12-B08-0159
%% lop: 12
%% chuong: 2
%% bai: 8
%% dang: 12.8.51
%% loai: TN4
%% muc: TH
%% dap_an: "C"
%% nguon: "CĐ2.4.pdf giáo viên gửi 27/10/2026 (M07, Lê Minh Khôi – Nguyễn Anh Khôi, Chương II Chủ đề 4 Các bài toán về tâm tỉ cự), Câu 9"
%% kien_thuc: "Bài 8 (biểu thức tọa độ của phép toán vectơ, độ dài vectơ), Bài 6 (điểm thỏa hệ thức vectơ, quy tắc ba điểm)"
%% trang_thai: chinh-thuc
%% ghi_chu: "Giáo viên đã duyệt 27/10/2026"
\begin{ex}
Trong không gian $Oxyz$, cho tứ diện $ABCD$ có tọa độ các đỉnh $A(0;1;0)$, $B(1;0;1)$, $C(5;-1;1)$ và $D(3;-3;2)$. Gọi $M(a;b;c)$ là điểm thỏa mãn $\left|\overrightarrow{MA}+\overrightarrow{MB}+\overrightarrow{MC}+\overrightarrow{MD}\right|$ đạt giá trị nhỏ nhất. Tính giá trị của biểu thức $R=a+b+c$.
\choice{$R=\dfrac32$}{$R=\dfrac12$}{\True $R=\dfrac52$}{$R=\dfrac72$}
\loigiai{$M$ là trọng tâm tứ diện: $M\left(\dfrac94;-\dfrac34;1\right)$. Vậy $R=\dfrac52$.}
\end{ex}
